\section{Описание практической части}
\label{sec:Chapter4} \index{Chapter4}

%Если в рамках работы писался какой-то код, здесь должно быть его
%описание: выбранный язык и библиотеки и мотивы выбора, архитектура,
%схема функционирования, теоретическая сложность алгоритма, характеристики
%функционирования (скорость/память).

Данная глава состоит из двух частей. В первой части рассматриваются ключевые аспекты реализации механизма репликации в исходном коде Linux 6.6. Во второй части приводятся результаты экспериментальной оценки производительности на тестовых приложениях.

\subsection{Детали реализации механизма в исходном коде}

\subsubsection{Виртуальный уровень и политики}

Одной из ключевых структур данных подсистемы управления памятью Linux является \texttt{vm\_area\_struct}. Она описывает диапазон пользовательской памяти, содержит информацию о его свойствах и семантике. Практически любая операция управления памятью в ядре начинается с обращения к соответствующему экземпляру \texttt{vm\_area\_struct}.

Одним из важнейших атрибутов данной структуры являются флаги (см листинг~\ref{lst:vm_area_struct}). Они представляют собой битовую маску, в которой каждый бит соответствует определённому свойству области памяти.

\begin{lstlisting}[
    style=normal,
    caption={Флаги виртуальной области памяти},
    label={lst:vm_area_struct}
]
/* include/linux/mm_types.h */

struct vm_area_struct {
    ...
	union {
		const vm_flags_t vm_flags;
		vm_flags_t __private __vm_flags;
	};
    ...
};
\end{lstlisting}

Для поддержки механизма репликации были добавлены два новых флага, отвечающих за репликацию таблиц трансляций и данных соответственно (см листинг~\ref{lst:vm_replica_flags}). Именно по ним и будет определяться поведение механизма репликации в отношении конкретной области памяти.

\begin{lstlisting}[style=diff]
 /* include/linux/mm.h */

 #ifdef CONFIG_ARCH_USES_HIGH_VMA_FLAGS
 #define VM_HIGH_ARCH_BIT_0	32
 ...
 #define VM_HIGH_ARCH_BIT_5	37
<+>+#define VM_HIGH_ARCH_BIT_6	39</+>
<+>+#define VM_HIGH_ARCH_BIT_7	40</+>
 #define VM_HIGH_ARCH_0	BIT(VM_HIGH_ARCH_BIT_0)
 ...
 #define VM_HIGH_ARCH_5	BIT(VM_HIGH_ARCH_BIT_5)
<+>+#define VM_HIGH_ARCH_6	BIT(VM_HIGH_ARCH_BIT_6)</+>
<+>+#define VM_HIGH_ARCH_7	BIT(VM_HIGH_ARCH_BIT_7)</+>
 #endif /* CONFIG_ARCH_USES_HIGH_VMA_FLAGS */
 ...
<+>+#ifdef CONFIG_USER_REPLICATION</+>
<+>+/* Page tables for this vma should be replicated during page faults */</+>
<+>+#define VM_REPLICA_INIT	VM_HIGH_ARCH_6</+>
<+>+/*</+>
<+>+ * Phys memory of this vma might/should be replicated.</+>
<+>+ * If this flag is set, VM_REPLICA_INIT also must be set.</+>
<+>+ */</+>
<+>+#define VM_REPLICA_COMMIT	VM_HIGH_ARCH_7</+>
<+>+#endif /* CONFIG_USER_REPLICATION */</+>
\end{lstlisting}
\begin{lstlisting}[
    style=normal,
    caption={Флаги репликации для виртуальных областей памяти},
    label={lst:vm_replica_flags}
]
/* include/linux/numa_user_replication.h */

#ifdef CONFIG_USER_REPLICATION
...
static inline bool numa_is_vma_replicant(struct vm_area_struct *vma) {
	return vma->vm_flags & VM_REPLICA_INIT;
}

static inline bool vma_has_replicas(struct vm_area_struct *vma) {
	return vma->vm_flags & VM_REPLICA_COMMIT;
}
...
#else
...
static inline bool numa_is_vma_replicant(struct vm_area_struct *vma) {
	return 0;
}

static inline bool vma_has_replicas(struct vm_area_struct *vma) {
	return 0;
}
...
#endif /* CONFIG_USER_REPLICATION */
\end{lstlisting}

Однако сами флаги репликации устанавливаются в соответствии с выбранными политиками репликации, на которых основан пользовательский интерфейс механизма (см главу~\ref{sec:Chapter3_user_API}). Политики задаются на уровне виртуального адресного пространства процесса, которое в Linux представляется структурой \texttt{mm\_struct}. Для хранения текущих настроек репликации в неё было добавлено соответствующее поле (см листинг~\ref{lst:mm_struct}).

\begin{lstlisting}[style=diff]
 /* include/linux/mm_types.h */

<+>+enum table_replication_policy_t {</+>
<+>+	/* no table replication */</+>
<+>+	TABLE_REPLICATION_NONE = 0,</+>
<+>+	/* replicate minimal amount to support RO-data replication */</+>
<+>+	TABLE_REPLICATION_MINIMAL,</+>
<+>+	/* replicate all tables */</+>
<+>+	TABLE_REPLICATION_ALL</+>
<+>+};</+>
<+>+</+>
<+>+enum data_replication_policy_t {</+>
<+>+	/* no data replicas at all */</+>
<+>+	DATA_REPLICATION_NONE = 0,</+>
<+>+	/* replicas are created on demand via numa balancer */</+>
<+>+	DATA_REPLICATION_ON_DEMAND,</+>
<+>+	/* all ro-data always replicated */</+>
<+>+	DATA_REPLICATION_ALL</+>
<+>+};</+>
<+>+</+>
<+>+struct numa_replication_policy_t {</+>
<+>+	enum table_replication_policy_t table_policy:3;</+>
<+>+	enum data_replication_policy_t data_policy:3;</+>
<+>+};</+>
<+>+</+>
 ...
 struct mm_struct {
     ...
<+>+#ifdef CONFIG_USER_REPLICATION</+>
<+>+    struct numa_replication_policy_t replication_policy;</+>
<+>+#endif</+>
 };
\end{lstlisting}

\begin{lstlisting}[
    style=normal,
    caption={Политики NUMA-репликации},
    label={lst:mm_struct}
]
/* include/linux/numa_user_replication.h */

#ifdef CONFIG_USER_REPLICATION
...
static inline enum table_replication_policy_t
get_table_replication_policy(struct mm_struct *mm) {
	return mm->replication_policy.table_policy;
}

static inline enum data_replication_policy_t
get_data_replication_policy(struct mm_struct *mm) {
	return mm->replication_policy.data_policy;
}
...
#endif /* CONFIG_USER_REPLICATION */
\end{lstlisting}

Таким образом, реализована инфраструктура, необходимая для назначения виртуальным областям памяти соответствующих флагов репликации. Установка флагов должна выполняться всякий раз, когда создаётся новая виртуальная область либо изменяются атрибуты уже существующей области, например, при обработке системных вызовов \texttt{mmap} и \texttt{mprotect}. Соответствующие изменения показаны в листинге~\ref{lst:vmflags_modify_replica}.

\begin{lstlisting}[style=normal]
/* include/linux/numa_user_replication.h */

#ifdef CONFIG_USER_REPLICATION
...
/* true if vma with @vm_flags is suitable for data replication */
static inline bool __vmflags_replica_candidate(vm_flags_t vm_flags) {
	return (vm_flags & VM_ACCESS_FLAGS) &&
        !(vm_flags & (VM_WRITE | VM_SHARED));
}

/* true if vma with @vm_flags is ready for VM_REPLICA_COMMIT flag, i.e.
 * it is suitable for data replication and VM_REPLICA_INIT is already
 * set */
static inline bool vmflags_might_be_replicated(vm_flags_t vm_flags) {
	return (vm_flags & VM_REPLICA_INIT) && (vm_flags & VM_ACCESS_FLAGS) &&
        !(vm_flags & (VM_WRITE | VM_SHARED));
}

/* set replication flags of @vm_flags according to @mm's policies */
static inline void vmflags_modify_replica(
    struct mm_struct *mm, unsigned long *vm_flags) {
	*vm_flags &= ~(VM_REPLICA_INIT | VM_REPLICA_COMMIT);

	if ((get_table_replication_policy(mm) == TABLE_REPLICATION_ALL) ||
		(get_table_replication_policy(mm) == TABLE_REPLICATION_MINIMAL &&
			__vmflags_replica_candidate(*vm_flags))) {
		*vm_flags |= VM_REPLICA_INIT;
		if (get_data_replication_policy(mm) != DATA_REPLICATION_NONE &&
				vmflags_might_be_replicated(*vm_flags))
			*vm_flags |= VM_REPLICA_COMMIT;
	}
}
...
#endif /* CONFIG_USER_REPLICATION */
\end{lstlisting}
\begin{lstlisting}[
    style=diff,
    caption={Установка флагов репликации новым/модифицированным областям},
    label={lst:vmflags_modify_replica}
]
 /* mm/mmap.c */

 unsigned long do_mmap(struct file *file, unsigned long addr,
             unsigned long len, unsigned long prot,
 	         unsigned long flags, vm_flags_t vm_flags,
 			 unsigned long pgoff, unsigned long *populate,
 			 struct list_head *uf)
 {
     struct mm_struct *mm = current->mm;
     ...
<+>+#ifdef CONFIG_USER_REPLICATION</+>
<+>+    vmflags_modify_replica(mm, &vm_flags);</+>
<+>+#endif</+>
     addr = mmap_region(file, addr, len, vm_flags, pgoff, uf);
     ...
     return addr;
 }

 /* mm/mprotect.c */

 static int do_mprotect_pkey(unsigned long start, size_t len,
         unsigned long prot, int pkey)
 {
     unsigned long end;
     struct vm_area_struct *vma;
     struct vma_iterator vmi;
     ...
     for_each_vma_range(vmi, vma, end) {
         unsigned long mask_off_old_flags, newflags;
         int new_vma_pkey;
         ...
         newflags = calc_vm_prot_bits(prot, new_vma_pkey);
         newflags |= (vma->vm_flags & ~mask_off_old_flags);
<+>+#ifdef CONFIG_USER_REPLICATION</+>
<+>+        if (vma->vm_mm)</+>
<+>+            vmflags_modify_replica(vma->vm_mm, &newflags);</+>
<+>+#endif</+>
         ...
     }
     ...
 }
\end{lstlisting}

Помимо этого флаги репликации приходится синхронно обновлять существующим областям при изменении процессом политик репликации таблиц трансляций и/или данных. На листинге~\ref{lst:numa_mm_handle_table_replication} показан код для смен политик репликации таблиц с \texttt{All} на \texttt{Minimal} и с \texttt{None} на \texttt{Minimal}. Кроме обновления флагов подходящим виртуальным областям инициируются операции репликации и дерепликации таблиц трансляций. Код для остальных переходов между политиками устроен аналогичным образом и отличается лишь набором областей, состояние которых необходимо изменить.

\begin{lstlisting}[
    style=normal,
    caption={Установка флага репликации таблиц трансляций при смене политики},
    label={lst:numa_mm_handle_table_replication}
]
/* mm/numa_user_replication.c */

/* called when table replication policy goes from All to Minimal */
static void dereplicate_all2min_pgtables(struct mm_struct *mm) {
	struct vm_area_struct *vma;
    ...
	do {
		unsigned long addr = vma->vm_start;
		struct vm_area_struct *next;
        ...
        /* clear VM_REPLICA_INIT and dereplicate
         * pgtables for not replica candidate vmas */
		if (numa_is_vma_replicant(vma) && !vma_might_be_replicated(vma)) {
			vm_flags_clear(vma, VM_REPLICA_INIT);
			...
			dereplicate_pgtables_pgd_range(&tlb, addr, vma->vm_end, ...);
		}
		vma = next;
	} while (vma);
    ...
}

static void numa_mm_table_replication_minimal(struct mm_struct *mm) {
	switch (get_table_replication_policy(mm)) {
    /* executed when table replication policy goes from None to Minimal */
	case TABLE_REPLICATION_NONE: {
		struct vm_area_struct *vma;
		MA_STATE(mas, &mm->mm_mt, 0, 0);

		mas_for_each(&mas, vma, ULONG_MAX) {
            /* set VM_REPLICA_INIT and replicate
             * pgtables for replica candidate vmas */
			if (vma_replica_candidate(vma)) {
				vm_flags_set(vma, VM_REPLICA_INIT);
				replicate_pgtables_vma(vma);
			}
		}
		break;
	}
	case TABLE_REPLICATION_MINIMAL:
		return;
	case TABLE_REPLICATION_ALL:
		dereplicate_all2min_pgtables(mm);
		break;
	}

	set_table_replication_policy(mm, TABLE_REPLICATION_MINIMAL);
}
\end{lstlisting}

Аналогичным образом реализовано изменение политики репликации данных. На листинге~\ref{lst:numa_mm_handle_data_replication} показаны переходы с \texttt{Lazy}/\texttt{All} на \texttt{None} и с \texttt{Lazy/None} на \texttt{All}. Вместе со сменой флагов выполняется репликация и дерепликация соответствующих физических страниц. Однако следует отметить, что при переходе с \texttt{None} на \texttt{Lazy} выполняется только обновление флагов, репликация данных осуществляется отложено при обработках \textit{NUMA hinting faults} (см главу~\ref{sec:Chapter3_data_repl_policies}).

\begin{lstlisting}[
    style=normal,
    caption={Установка флага репликации данных при смене политики},
    label={lst:numa_mm_handle_data_replication}
]
/* mm/numa_user_replication.c */

static bool numa_mm_data_replication_none(struct mm_struct *mm) {
	switch (get_data_replication_policy(mm)) {
	case DATA_REPLICATION_NONE:
		return true;
    /* executed when data replication policy
     * goes from Lazy/All to None */
	case DATA_REPLICATION_ON_DEMAND:
	case DATA_REPLICATION_ALL: {
		struct vm_area_struct *vma;
		MA_STATE(mas, &mm->mm_mt, 0, 0);

		mas_for_each(&mas, vma, ULONG_MAX) {
            /* deduplicate data and clear VM_REPLICA_COMMIT
             * for all has_replicas vmas */
			if (vma_has_replicas(vma)) {
				vma_start_write(vma);
				if (phys_deduplicate(vma,
						vma->vm_start, vma->vm_end - vma->vm_start))
					return false;
				__vm_flags_mod(vma, VM_NONE, VM_REPLICA_COMMIT);
			}
		}
		break;
	}
	}

	set_data_replication_policy(mm, DATA_REPLICATION_NONE);
	return true;
}

static bool numa_mm_data_replication_all(struct mm_struct *mm) {
	switch (get_data_replication_policy(mm)) {
    /* executed when data replication policy
     * goes from Lazy/None to All */
	case DATA_REPLICATION_NONE:
	case DATA_REPLICATION_ON_DEMAND: {
		struct vm_area_struct *vma;
		MA_STATE(mas, &mm->mm_mt, 0, 0);

		mas_for_each(&mas, vma, ULONG_MAX) {
            /* set VM_REPLICA_COMMIT an duplicate data
             * for all replica candidate vmas */
			if (vma_might_be_replicated(vma))
				vm_flags_set(vma, VM_REPLICA_COMMIT);
			if (vma_has_replicas(vma))
				phys_duplicate(vma,
                    vma->vm_start, vma->vm_end - vma->vm_start);
		}
		break;
	}
	case DATA_REPLICATION_ALL:
		return true;
	}

	set_data_replication_policy(mm, DATA);
	return true;
}
\end{lstlisting}

\subsubsection{Физический уровень}
\label{sec:Chapter4_phys_layer}

Перед тем как перейти непосредственно к репликации/дерепликации физических страниц, будь то страницы с таблицами трансляций или страницы с данными, необходимо рассмотреть несколько вспомогательных сущностей физического уровня, введённых для нужд репликации.

Физическое адресное пространство разбито на страницы, каждая из которых в Linux имеет \texttt{page}-дескриптор~--- экземпляр структуры \texttt{page}, содержащий всю информацию о ней. Как и в случае виртуальных областей памяти, одним из основных атрибутов физической страницы, являются флаги (см листинг~\ref{lst:page})~--- они также представляют собой битовую маску, в которой каждый бит соответствует определённому свойству страницы.

\begin{lstlisting}[
    style=normal,
    caption={Флаги физической страницы},
    label={lst:page}
]
/* include/linux/mm_types.h */

struct page {
    unsigned long flags;
    ...
};
\end{lstlisting}

Для обозначения реплицированности физической страницы был добавлен флаг \texttt{PG\_replicated} (см листинг~\ref{lst:PG_replicated}). Он выставляется у всех реплицированных страниц~--- как у страниц с пользовательскими данными, так и у страниц с таблицами трансляций.

\begin{lstlisting}[
	style=diff,
	caption={Флаг для реплицированных физических страниц},
	label={lst:PG_replicated}
]
 /* include/linux/page-flags.h */

 enum pageflags {
	 PG_locked,
	 ...
<+>+</+>
<+>+#ifdef CONFIG_USER_REPLICATION</+>
<+>+	 PG_replicated,</+>
<+>+#endif</+>
 	 __NR_PAGEFLAGS,
	 ...
 };
 ...
<+>+#ifdef CONFIG_USER_REPLICATION</+>
<+>+PAGEFLAG(Replicated, replicated, PF_ANY)</+>
<+>+#else</+>
<+>+PAGEFLAG_FALSE(Replicated, replicated)</+>
<+>+#endif</+>
<+>+</+>
 ...
<+>+#ifdef CONFIG_USER_REPLICATION</+>
<+>+#define __PG_REPLICATED		(1UL << PG_replicated)</+>
<+>+#else</+>
<+>+#define __PG_REPLICATED		0</+>
<+>+#endif</+>
<+>+</+>
 /*
  * Flags checked when a page is freed.  Pages being freed should not have
  * these flags set.  If they are, there is a problem.
  */
 #define PAGE_FLAGS_CHECK_AT_FREE					\
 	 1UL << PG_private    | 1UL << PG_private_2 |	\
 	 1UL << PG_writeback  | 1UL << PG_reserved  |	\
 	 1UL << PG_slab       | 1UL << PG_active    |	\
<->-	 1UL << PG_unevictable| __PG_MLOCKED | LRU_GEN_MASK)</->
<+>+	 1UL << PG_unevictable| __PG_REPLICATED | __PG_MLOCKED | LRU_GEN_MASK)</+>
\end{lstlisting}

Такое архитектурное решение накладывает небольшое ограничение на репликацию данных, которые могут разделяться между несколькими процессами. Необходимо исключать ситуации, когда одна и та же страница с данными реплицирована в контексте одного процесса и не реплицирована в контексте другого (см рис.~\ref{fig:data_repl_err}), поскольку значение флага \texttt{PG\_replicated} является общим для страницы. Именно поэтому при репликации и дерепликации данных исходная страница никогда не переиспользуется в качестве страницы-реплики (см главу~\ref{sec:Chapter3_data_repl}), а страница-реплика~--- в качестве оригинальной страницы (см главу~\ref{sec:Chapter3_data_table_derepl}).

Рассмотрим теперь особенности реализации реплицированных таблиц трансляций. В Linux для таблиц уровней PMD и PTE ассоциированы собственные блокировки для сериализации их модификаций. Чтобы согласованно модифицировать таблицы-реплики, необходимо использовать общую для них блокировку~--- об этом уже шла речь в предыдущих главах. Для этого среди таблиц-реплик фиксируется одна \textit{мастер}-таблица, блокировка которой и используется. И чтобы иметь доступ к этой блокировке из контекста любой таблицы-реплики, в каждой из них хранится указатель на \textit{мастер}-таблицу.

\begin{figure}[H]
    \centering
    \includegraphics[width=1\textwidth]{data_repl_err.pdf}
    \caption{Невалидное состояние страницы с данными, при котором для одного процесса она реплицирована, а для другого~--- нет}
    \label{fig:data_repl_err}
\end{figure}

Для указателя на \textit{мастер}-таблицу добавлено специальное поле в \texttt{ptdesc}-дескрипторе~--- структуре-алиасе \texttt{page}-дескриптора, содержащей поля специально для таблиц трансляций (см листинг~\ref{lst:master_table}). Новое поле добавлено в одно из существующих объединений, поэтому размер структуры остаётся неизменным, а с другими полями объединения оно не конфликтует, т.к. они используются только в архитектурах x86 и PowerPC, не в arm64. Первоначально каждая таблица является сама себе \textit{мастер}-таблицей, однако в процессе репликации соответствующие указатели обновляются.

\begin{lstlisting}[
	style=diff,
	caption={Использование блокировки \textit{мастер}-таблицы},
	label={lst:master_table}
]
 /* include/linux/mm_types.h */

 struct ptdesc {
	 ...
	 union {
		 struct mm_struct *pt_mm; /* Used for x86 pgds. */
		 atomic_t pt_frag_refcount; /* Powerpc only. */
<+>+#ifdef CONFIG_USER_REPLICATION</+>
<+>+		 struct ptdesc *master_table;</+>
<+>+#endif</+>
	 };

	 union {
 #if ALLOC_SPLIT_PTLOCKS
		 spinlock_t *ptl;
 #else
		 spinlock_t ptl;
 #endif
	 };
	 ...
 };

 /* include/linux/mm.h */

 #if USE_SPLIT_PTE_PTLOCKS
 #if ALLOC_SPLIT_PTLOCKS
 ...
 bool ptlock_alloc(struct ptdesc *ptdesc);
 void ptlock_free(struct ptdesc *ptdesc);

 static inline spinlock_t *ptlock_ptr(struct ptdesc *ptdesc) {
<+>+#ifdef CONFIG_USER_REPLICATION</+>
<+>+	 return ptdesc->master_table->ptl;</+>
<+>+#else</+>
	 return ptdesc->ptl;
<+>+#endif</+>
 }
 #else /* ALLOC_SPLIT_PTLOCKS */
 ...
 static inline bool ptlock_alloc(struct ptdesc *ptdesc) {
<+>+#ifdef CONFIG_USER_REPLICATION</+>
<+>+	 ptdesc->master_table = ptdesc;</+>
<+>+#endif</+>
	 return true;
 }

 static inline void ptlock_free(struct ptdesc *ptdesc) {
<+>+#ifdef CONFIG_USER_REPLICATION</+>
<+>+	 ptdesc->master_table = NULL;</+>
<+>+#endif</+>
 }

 static inline spinlock_t *ptlock_ptr(struct ptdesc *ptdesc) {
<+>+#ifdef CONFIG_USER_REPLICATION</+>
<+>+	 return &ptdesc->master_table->ptl;</+>
<+>+#else</+>
	 return &ptdesc->ptl;
<+>+#endif</+>
 }
 #endif /* ALLOC_SPLIT_PTLOCKS */
 ...
 #endif /* USE_SPLIT_PTE_PTLOCKS */
\end{lstlisting}

Следующая вспомогательная сущность тоже касается реплицированных таблиц трансляций. Для согласованного изменения всех таблиц-реплик необходимо иметь возможность последовательно обходить их. Для этого они объединяются в циклический односвязный список с помощью нового поля в \texttt{ptdesc}-дескрипторе (см листинг~\ref{lst:replica_list}). Размер структуры всё также остаётся неизменным, потому что новое поле добавлено в объединение с уже существующим. И конфликтов с ним также не возникнет, поскольку существующее поле является алиасом одноимённого поля структуры \texttt{page}, не используемого для таблиц трансляций. В процессе репликации таблицы-реплики соединяются в список, после чего его обход может выполняться с помощью макросов \texttt{for\_each\_pgtable} и \texttt{for\_each\_pgtable\_replica} (см листинг~\ref{lst:replica_list}).

\begin{lstlisting}[style=diff]
 /* include/linux/llist.h */

 struct llist_head {
	 struct llist_node *first;
 };
 
 struct llist_node {
	 struct llist_node *next;
 };

 /* include/linux/mm_types.h */

 struct ptdesc {
	 ...
<->-	 unsigned long __page_mapping; /* Unused for page tables. */</->
<+>+	 union {</+>
<+>+		 unsigned long __page_mapping; /* Unused for page tables. */</+>
<+>+		 struct llist_head replica_list_head;</+>
<+>+		 struct llist_node replica_list_node;</+>
<+>+	 };</+>
	 ...
 };
\end{lstlisting}
\begin{lstlisting}[
	style=normal,
	caption={Циклический односвязный список таблиц-реплик},
	label={lst:replica_list}
]
/* include/linux/numa_user_replication.h */

#ifdef CONFIG_USER_REPLICATION
...
/**
 * @pos:	struct ptdesc* of current replica
 * @table:	current table entry to write (virtual address)
 * @head_table:	table entry to start with, will be a part of this loop
 * @nid:	node id of current pgtable
 * @offset:	offset of current table entry in table page in bytes
 * @start:	boolean value for tmp storage
 */
#define for_each_pgtable(pos, table, head_table, nid, offset, start)	\
	for (pos = llist_entry(									\
			&virt_to_ptdesc(head_table)->replica_list_node,	\
		typeof(*pos), replica_list_node),					\
	    start = true, nid = page_to_nid(ptdesc_page(pos)),	\
	    offset = offset_in_table(head_table),				\
		table = (typeof(table))get_table_ptr(pos, offset);	\
	    pos != virt_to_ptdesc(head_table) || start;			\
	    pos = llist_entry((pos)->replica_list_node.next,	\
			typeof(*pos), replica_list_node),				\
	    table = (typeof(table))get_table_ptr(pos, offset),	\
	    nid = page_to_nid(ptdesc_page(pos)), start = false)

/**
 * @pos:	struct ptdesc* of current replica
 * @table:	current table entry to write (virtual address)
 * @head_table:	table entry to start with, will not be a part of this loop
 * @offset:	offset of current table entry in table page in bytes
 */
#define for_each_pgtable_replica(pos, table, head_table, offset)	\
	for (pos = llist_entry(										\
			virt_to_ptdesc(head_table)->replica_list_node.next,	\
		typeof(*pos), replica_list_node),					\
	    offset = offset_in_table(head_table),				\
		table = (typeof(table))get_table_ptr(pos, offset);	\
	    pos != virt_to_ptdesc(head_table);					\
	    pos = llist_entry((pos)->replica_list_node.next,	\
			typeof(*pos), replica_list_node),				\
	    table = (typeof(table))get_table_ptr(pos, offset))
...
#endif /* CONFIG_USER_REPLICATION */
\end{lstlisting}

Последнее, о чём в данной главе стоит сказать --- это корень иерархии таблиц трансляций, PGD-уровень. В отсутствии репликации он состоит из единственной PGD-таблицы, указатель на которую хранится в специальном поле \texttt{pgd} структуры \texttt{mm\_struct}. Но при использовании механизма репликации для каждой NUMA-ноды создаётся собственная PGD-таблица. Такие таблицы аллоцируются при иницилизации нового экземпляра \texttt{mm\_struct}, а указатели на них сохраняются в специальном массиве, доступном через новое поле данной структуры (см листинг~\ref{lst:pgd_numa}).

\begin{lstlisting}[style=diff]
 /* include/linux/mm_types.h */

 struct mm_struct {
	 struct {
		 ...
		 pgd_t *pgd;
		 ...
	 };
<+>+#ifdef CONFIG_USER_REPLICATION</+>
<+>+	 pgd_t **pgd_numa;</+>
<+>+#endif</+>
	 ...
 };
\end{lstlisting}
\begin{lstlisting}[
	style=normal,
	caption={Указатели на реплицированные PGD-таблицы доступны через \texttt{mm\_struct}},
	label={lst:pgd_numa}
]
/* include/linux/numa_user_replication.h */

#ifdef CONFIG_USER_REPLICATION
...
#define this_node_pgd(mm) ((mm)->pgd_numa[numa_node_id()])
...
#else
...
#define this_node_pgd(mm) ((mm)->pgd)
...
#endif /* CONFIG_USER_REPLICATION */
\end{lstlisting}

% Моя стратегия теперь следующая:
% 1. Сначала рассказать про виртуальный уровень. Это то, с чего и нужно начинать, потому что это плавный переход от интерфейса политик, к недрам реализации.
% 2. Во второй подглаве я пожалуй сделаю отступление и расскажу про тонкости физического уровня:
%    * PG_replicated флаг;
%    * master page table lock (единый для всех таблиц-реплик);
%    * кольцевой список таблиц-реплик для итерирования;
%    * массив pgd_numa со всеми PGD_таблицами-репликами.
% 3. Далее в ход пойдет репликация таблиц на page fault-ах.
% 4. Репликация даных + модель ее вызова при политиках Lazy и All.
% 5. Дерепликация таблиц и данных.
